% Plant AE Paper — System Characterisation and Experimental Validation
% CB/JH, THA Augsburg, Projekt I1, 2026
\documentclass[10pt,a4paper,twocolumn]{article}
\usepackage[top=2.0cm,bottom=2.2cm,left=1.6cm,right=1.6cm,columnsep=0.6cm]{geometry}
\usepackage{amsmath,amssymb}
\usepackage{graphicx}
\usepackage{booktabs}
\usepackage{siunitx}
\usepackage[hidelinks]{hyperref}
\usepackage{microtype}
\usepackage{xcolor}
\usepackage{caption}
\usepackage{enumitem}
\usepackage{fancyhdr}
\usepackage{titlesec}
\usepackage[round,authoryear]{natbib}
\usepackage{placeins}

\hypersetup{colorlinks=true,linkcolor=blue!60!black,citecolor=green!50!black}
\captionsetup{font=small,labelfont=bf,skip=4pt}
\setlength{\abovecaptionskip}{4pt}
\setlength{\belowcaptionskip}{0pt}
\titleformat{\section}{\normalfont\bfseries\normalsize}{\thesection.}{0.5em}{}
\titleformat{\subsection}{\normalfont\itshape\normalsize}{\thesubsection}{0.5em}{}
\titlespacing{\section}{0pt}{6pt}{3pt}
\titlespacing{\subsection}{0pt}{4pt}{2pt}
\setlength{\parskip}{2pt}

% Float placement: allow more figure area per page so wide figures sit
% near their reference instead of drifting pages later. Full-width
% (figure*) floats obey the \dbl... parameters in two-column mode.
\renewcommand{\topfraction}{0.92}
\renewcommand{\bottomfraction}{0.85}
\renewcommand{\textfraction}{0.07}
\renewcommand{\floatpagefraction}{0.72}
\setcounter{topnumber}{3}
\setcounter{totalnumber}{4}
\renewcommand{\dbltopfraction}{0.92}
\renewcommand{\dblfloatpagefraction}{0.72}
\setcounter{dbltopnumber}{3}

\newcommand{\gammasq}{\ensuremath{\hat{\gamma}^2}}
\newcommand{\taumle}{\ensuremath{\hat{\tau}_{\mathrm{MLE}}}}

\pagestyle{fancy}\fancyhf{}
\renewcommand{\headrulewidth}{0.4pt}
\lhead{\small\textit{Plant AE Monitoring: System Characterisation}}
\rhead{\small\textit{CB/JH, THA Augsburg, Projekt I1, 2026}}
\cfoot{\small\thepage}

\begin{document}

\twocolumn[{%
\begin{@twocolumnfalse}
\begin{center}
{\LARGE\bfseries Passive Acoustic Emission Monitoring of Plant Stress Responses:\\[3pt]
System Characterisation and Experimental Validation}\\[8pt]
{\normalsize Christoph Bellmann$^{1}$, Julian Hübke$^{1}$}\\[3pt]
{\small $^{1}$THA Augsburg, Systems Engineering, Projekt I1, Sommersemester 2026}\\[2pt]
{\small \texttt{julian.huebke@tha.de, christoph.bellmann@tha.de}}\\[8pt]
\end{center}
\noindent\textbf{Abstract.} We present the design, characterisation, and experimental evaluation of a low-cost passive acoustic emission (AE) monitoring system for plant physiological stress detection. The original system uses three piezoelectric channels; a Phase-3 follow-up replaces CH3/CH4 with MEMS microphones. All channels are amplified by LM358 circuits and recorded at \SI{500}{\kilo\SIUnitSymbolMicro s} in deep-memory mode (\num{300000} samples, \SI{0.6}{\second} per frame) on a Rigol DS1104Z. Frequency characterisation via a Raspberry Pi Pico chirp sweep identifies a calibrated piezo band of \SIrange{5}{35}{\kilo\hertz} (multi-frame coherence $\gammasq \geq 0.95$). Multi-frame MLE time-delay estimation confirms equidistant sensor placement ($|\taumle|<\SI{0.3}{\micro\second}$); a $+\SI{1.95}{\micro\second}$ offset at the \SI{9.17}{\kilo\hertz} steel eigenmode is a resonance-phase artefact. A controlled watering experiment yields no detectable AE after drift correction (Mann--Whitney $p=0.86$). In the MEMS follow-up, moving both microphones from air into soil raises median RMS from \SI{80.0}{\milli\volt} to \SI{302.8}{\milli\volt}, yet both channels remain nearly simultaneous (GCC-PHAT median \SI{0.000}{\micro\second}) and share a \SI{32.59}{\kilo\hertz} line with $\gammasq=0.997$. No first--second-half band comparison survives FDR correction (0/32). The combined evidence indicates common-mode interference and one CH3 saturation fault, not plant AE. A positive-control reference experiment confirms pipeline integrity (\SI{+52.6}{\deci\bel} sensitivity at \SI{290}{\hertz}, Wilcoxon $p<10^{-12}$).\par\bigskip
\end{@twocolumnfalse}
}]

%% ─────────────────────────────────────────────────────────────────────────────
\section{Introduction}
%% ─────────────────────────────────────────────────────────────────────────────

Acoustic emission (AE) in plants (transient elastic waves generated by rapid internal stress redistributions) has been linked to xylem cavitation, cell-wall micro-fracture, and turgor-driven growth \citep{Milburn1966,Tyree1983,Zweifel2012}. Detecting these signals non-invasively would provide a real-time proxy for plant water stress, relevant for precision irrigation and phenotyping. Classical AE plant research employs resonant ultrasonic sensors above \SI{100}{\kilo\hertz} \citep{Cochard1992,Nardini2015,DeRoo2016}. More recent work extends AE sensing to airborne and contactless detection \citep{Khait2023,Bonisoli2025,Klaminder2025}, positions it as an early stress indicator for precision agriculture \citep{Oletic2020,Kiraly2025,Saha2026}, and demonstrates first AE-driven irrigation prototypes \citep{Oletic2023,Simbeye2023}. Low-cost alternatives using standard op-amps and off-the-shelf oscilloscopes are attractive for continuous monitoring but their sensitivity limits are rarely rigorously quantified.

This work addresses that gap with five contributions: (i)~full frequency-domain characterisation of a three-channel LM358-based acquisition chain; (ii)~multi-frame MLE time-delay estimation for sensor placement verification; (iii)~a statistically rigorous pre--post watering experiment with explicit drift correction; (iv)~identification of the dominant noise sources limiting plant AE detectability in this class of hardware; and (v)~a Phase-3 air-versus-soil evaluation of two MEMS channels using the same spectral, coherence, transient-screening, TDE, and FDR decision gates.

%% ─────────────────────────────────────────────────────────────────────────────
\section{Nomenclature}
%% ─────────────────────────────────────────────────────────────────────────────

\noindent\textit{\small Abbreviations.}
{\small
\begin{description}[leftmargin=1.55cm,labelwidth=1.3cm,labelsep=0.25cm,%
  labelindent=0pt,style=sameline,font=\normalfont,itemsep=0.6pt,topsep=2pt,parsep=0pt]
\item[AE] Acoustic emission
\item[AC / DC] Alternating / direct current (input coupling / offset)
\item[ADC] Analog-to-digital converter
\item[DFT] Discrete Fourier transform
\item[FDR] False discovery rate (Benjamini--Hochberg)
\item[GBW] Gain--bandwidth product (of the amplifier)
\item[MEMS] Micro-electro-mechanical system (microphone)
\item[MLE] Maximum-likelihood estimation
\item[PSD] Power spectral density (Welch estimate)
\item[PZT] Lead zirconate titanate (piezo material)
\item[RMS] Root mean square
\item[SNR] Signal-to-noise ratio
\item[TDE] Time-delay estimation
\item[TTL] Transistor--transistor logic (sync pulse)
\end{description}}

\smallskip
\noindent\textit{\small Symbols.}
{\small
\begin{description}[leftmargin=1.25cm,labelwidth=1.0cm,labelsep=0.25cm,%
  labelindent=0pt,style=sameline,font=\normalfont,itemsep=0.6pt,topsep=2pt,parsep=0pt]
\item[$\gammasq$] Magnitude-squared coherence between two channels ($0$--$1$; $1$ = identical waveform)
\item[$\taumle$] Maximum-likelihood time-delay estimate (narrowband phase; Eq.~\ref{eq:tde})
\item[$\hat\tau_\mathrm{XC}$] Cross-correlation time-delay estimate (cross-check)
\item[$\hat\tau_\mathrm{ph}$] Phase-difference time-delay estimate
\item[$E_b$] Mean band energy in a \SI{5}{\kilo\hertz} bin $b$
\item[$V_3,\,V_4$] Signal of rod channels 3 and 4 (DFT $X_n$; RMS in Tab.~\ref{tab:cal})
\item[$f_0$] Narrowband analysis frequency
\item[$f_{-3}$] \SI{-3}{\deci\bel} (high-pass) corner frequency
\item[$e_n$] Input-referred amplifier voltage-noise density (nV$/\sqrt{\text{Hz}}$)
\item[$N$] Number of frames
\item[$M$] Number of DFT bins within a band
\item[$\alpha$] Significance level (here $0.05$)
\item[$\tau$] Kendall rank correlation (trend); \emph{not} the time delay
\item[$p$] Significance probability ($p$-value)
\item[$\Delta E$] Band-energy change, post $-$ pre watering
\end{description}}

%% ─────────────────────────────────────────────────────────────────────────────
\section{Materials and Methods}
%% ─────────────────────────────────────────────────────────────────────────────

\subsection{Analysis Methods Overview}

The AE characterisation pipeline combines four method families, applied in sequence from raw frames to statistical decision. They are detailed in the subsections that follow and summarised here for orientation.

\begin{itemize}[leftmargin=1.2em,itemsep=1pt,topsep=2pt]
  \item \textbf{Pre-processing.} Per-frame DC removal, linear detrending, Hann windowing, and a clipping/plausibility check reject saturated captures before analysis.
  \item \textbf{Spectral characterisation.} Averaged FFT and Welch PSD ($\text{nperseg}=8192$) per channel; band energies $E_b$ in \SI{5}{\kilo\hertz} bins (\S\ref{sec:daq}); prominence-based peak detection with mains-harmonic rejection, extracting frequency, prominence, bandwidth, $Q$-factor, and SNR per peak.
  \item \textbf{AE detection and validation.} Baseline-versus-active peak comparison (only peaks emerging under excitation are retained), repeatability across repeated captures, cross-channel consistency, and impulse-response resonance screening to separate sensor eigenmodes from genuine emissions.
  \item \textbf{Inter-channel and temporal analysis.} Multi-frame coherence $\gammasq$ (Eq.~\ref{eq:coh}) and time-delay estimation by narrowband phase MLE with a cross-correlation cross-check (Eq.~\ref{eq:tde}).
  \item \textbf{Statistical testing.} Wilcoxon signed-rank (paired pre--post), Mann--Whitney $U$ (drift-corrected), Kendall $\tau$ (trend), all with Benjamini--Hochberg FDR control (\S\ref{sec:stats}).
\end{itemize}

\subsection{Hardware Setup}
\label{sec:hw}

Figure~\ref{fig:system} shows the measurement chain. Channel~1 (CH1) is a bare soil piezo disc (\SI{27}{\milli\metre}) inserted into the root zone, biased through \SI{830}{\kilo\ohm} and buffered by an LM358 amplifier. Channels~3 and~4 (CH3, CH4) use identical discs clamped to the ends of a \SI{60}{\milli\metre} steel rod via \SI{820}{\kilo\ohm} resistors and separate LM358 circuits; they serve as structural references. Channel~2 is defective and excluded. All channels use Rigol RP2200 passive probes at the $10\times$ position. Probe compensation was performed manually against the oscilloscope's \SI{1}{\kilo\hertz} calibration square wave; oscilloscope self-calibration was not possible because the defective CH2 prevents freeing all inputs simultaneously.

The setup was subsequently relocated to another place and plant: CH1 remains a piezo sensor connected through its existing LM amplifier, whereas the piezo sensors on CH3 and CH4 were replaced by MEMS microphones connected through the existing LM amplifier stages. The monitored specimen is a lettuce plant (\textit{Lactuca sativa}) rather than the \textit{Euphorbia leuconeura} used previously. Sections~3--5 retain the historical three-piezo results on \textit{Euphorbia leuconeura}; Section~\ref{sec:mems} reports the separately characterised MEMS air-reference and soil-insertion runs on \textit{Lactuca sativa}.

\begin{figure*}[!tb]
  \centering
  \includegraphics[width=0.82\textwidth]{figures/fig1_system.png}
  \caption{Main measurement system used to acquire the data analysed in this paper. Three piezo channels amplified by LM358 circuits are recorded by a Rigol DS1104Z at \SI{500}{\kilo Sa\per s}, \num{300000}~pts per frame.}
  \label{fig:system}
\end{figure*}

\begin{figure*}[!tb]
  \centering
  \includegraphics[width=0.82\textwidth]{figures/fig2_bandwidth.png}
  \caption{How faithfully the chain passes each frequency. A Raspberry Pi Pico injects a known \SIrange{20}{100}{\kilo\hertz} chirp; the two rod sensors (CH3, CH4) record it. (a)~Coherence $\gammasq$ between the two channels per \SI{5}{\kilo\hertz} band ($1$ = identical waveform, i.e.\ the frequency is passed cleanly); green shading marks the calibrated band. (b)~Recorded RMS amplitude of each channel; the peaks at \SI{9.17}{\kilo\hertz}, \SI{20}{\kilo\hertz}, \SI{30}{\kilo\hertz} are steel-rod eigenmodes.}
  \label{fig:bandwidth}
\end{figure*}

\begin{table}[!h]
\centering
\caption{Calibration band summary (Pico chirp sweep). Band in \si{\kilo\hertz}; $V_3,V_4$ are per-band RMS amplitudes (\si{\milli\volt}). Reliable: $\gammasq \geq 0.9$.}
\label{tab:cal}
\small
\setlength{\tabcolsep}{4pt}
\begin{tabular}{lrrrr}
\toprule
Band & $\gammasq$ & $V_3$ & $V_4$ & $\taumle$~(\si{\micro\second}) \\
\midrule
5  & 0.962 & 6.44  & 6.42  & -0.26 \\
10 & 1.000 & 105.7 & 98.6  & +0.19 \\
15 & 0.993 & 11.3  & 10.6  & -0.20 \\
20 & 0.976 & 32.6  & 31.0  & -0.06 \\
25 & 0.983 & 10.8  & 10.2  & -0.67 \\
30 & 0.996 & 53.4  & 51.7  & -0.08 \\
35 & 0.959 & 14.5  & 14.2  & +0.02 \\
40 & 0.064 & 4.15  & 3.92  & +4.29 \\
50 & 0.065 & 3.94  & 3.01  & +2.11 \\
\bottomrule
\end{tabular}
\end{table}


\subsection{Data Acquisition and Processing}
\label{sec:daq}

Measuring at \SI{500}{\kilo Sa\per s}, \num{300000}~samples (\SI{0.6}{\second}), AC coupling, \SI{0.5}{\volt/div}. A Raspberry Pi Pico generates a \SIrange{20}{100}{\kilo\hertz} chirp (\SI{30}{\second} sweep) for calibration, synchronised to the oscilloscope via a \SI{500}{\micro\second} TTL pulse at every \SI{5}{\kilo\hertz} milestone.

Each acquisition automatically commits its raw waveforms (\texttt{.npz}) to the project repository via \mbox{git-LFS}, alongside a per-run JSON/Markdown summary and metadata. This yields a versioned, append-only raw-data pipeline in which every reported figure is traceable to an immutable, content-addressed recording.

Band energies: $E_b = \frac{1}{NM}\sum_{n=1}^N\sum_{m\in b}|X_n(m)|^2$, where $X_n$ is the DFT of frame $n$ and $b$ is a \SI{5}{\kilo\hertz}-wide bin. Multi-frame coherence:
\begin{equation}
  \gammasq(f) = \frac{\bigl|\sum_n V_3(n,f)\,V_4^*(n,f)\bigr|^2}
                     {\sum_n|V_3|^2\cdot\sum_n|V_4|^2}.
  \label{eq:coh}
\end{equation}

\subsection{Time-Delay Estimation}

MLE at narrowband frequency $f_0$:
\begin{equation}
  \taumle = -\frac{\angle\bigl(\sum_n V_3(n,f_0)\,V_4^*(n,f_0)\bigr)}{2\pi f_0},
  \label{eq:tde}
\end{equation}
unambiguous for $|\tau|<(2f_0)^{-1}$. A cross-correlation estimate $\hat\tau_\text{XC}$ provides a cross-check.

\subsection{Statistical Methods}
\label{sec:stats}

Pre--post comparisons: Wilcoxon signed-rank test on paired frames. Multiple-band comparisons: Benjamini--Hochberg FDR at $\alpha=0.05$ \citep{Benjamini1995}. Trend detection: Kendall $\tau$. Drift-corrected test: Mann--Whitney $U$ comparing pre-watering steady-state (frames~0--4) against post-watering steady-state (frames~10--19).

%% ─────────────────────────────────────────────────────────────────────────────
\section{System Characterisation}
%% ─────────────────────────────────────────────────────────────────────────────

\subsection{Amplifier Bandwidth and Coherence}

Figure~\ref{fig:bandwidth} shows results from the Pico chirp sweep. Three regions are identified:

\begin{itemize}[leftmargin=1.2em,itemsep=1pt,topsep=2pt]
  \item \textbf{\SIrange{5}{35}{\kilo\hertz}}: $\gammasq \geq 0.95$, flat amplitude except at resonance peaks (\SI{9.17}{\kilo\hertz}, \SI{20}{\kilo\hertz}, \SI{30}{\kilo\hertz} steel eigenmodes). \emph{Calibrated measurement band.}
  \item \textbf{\SIrange{40}{55}{\kilo\hertz}}: $\gammasq = 0.06$--$0.45$; LM358 roll-off. 
  \item \textbf{\SIrange{60}{75}{\kilo\hertz}}: partial recovery ($\gammasq\approx 0.85$), tentatively spurious probe-cable resonance; unreliable.
\end{itemize}

The LM358 with \SI{830}{\kilo\ohm} parallel input resistance and estimated piezo capacitance \SI{10}{\nano\farad} yields high-pass corner $f_{-3}\approx\SI{19}{\hertz}$, leaving all mains harmonics $\geq\SI{50}{\hertz}$ unattenuated.

\subsection{Time-Delay Estimation}

Figure~\ref{fig:tde}a shows $\taumle$ for all 20 calibration bands. Within the calibrated band, all delays lie within $\pm\SI{0.3}{\micro\second}$, confirming equidistant sensor placement on the rod. Unreliable bands ($\gammasq<0.5$) show scatter up to $\pm\SI{4.3}{\micro\second}$, demonstrating that low coherence directly inflates TDE variance.

At the \SI{9.17}{\kilo\hertz} eigenmode, fine-resolution phase-difference estimation across $N=15$ clean frames yields $\hat\tau_\text{ph} = +\SI{1.950\pm0.031}{\micro\second}$ and $\hat\tau_\text{XC}=+\SI{1.60}{\micro\second}$ (Fig.~\ref{fig:tde}b). Both are consistent and physically interpretable as a resonance-induced group-delay offset \citep{Knapp1976}: at the eigenfrequency the rod acts as a delay line for the propagating mode, introducing a geometry-independent phase difference between sensor positions.

\begin{figure*}[!tb]
  \centering
  \includegraphics[width=0.82\textwidth]{figures/fig3_tde.png}
  \caption{Time-delay between the two rod sensors. (a)~MLE delay $\taumle$ per \SI{5}{\kilo\hertz} band; green circles are trustworthy ($\gammasq\geq0.9$), grey crosses are low-coherence bands whose scatter is meaningless. In the calibrated band all delays sit inside $\pm\SI{0.5}{\micro\second}$ (green band), i.e.\ the sensors are equidistant. (b)~Zoom on the \SI{9.17}{\kilo\hertz} steel resonance ($N=15$ clean frames): every frame gives the same $\approx\SI{1.95}{\micro\second}$ offset, a fixed resonance-phase artefact rather than a geometric path difference. Magnitudes shown; the phase and cross-correlation estimates agree to within \SI{0.35}{\micro\second}.}
  \label{fig:tde}
\end{figure*}

\subsection{Channel Noise Floor}

Passive band-energy spectra from the pre-watering baseline ($N=10$ frames) reveal three distinct channel noise regimes.

\textbf{CH4} (rod piezo~B) shows a flat spectrum of $\approx\SI{0.8}{\milli\volt}$ RMS per \SI{5}{\kilo\hertz} band, totalling \SI{25.7}{\milli\volt} RMS. This matches the LM358 noise prediction ($40\,\text{nV}/\sqrt{\text{Hz}} \times \sqrt{500\,\text{kHz}} \approx \SI{28}{\milli\volt}$), confirming CH4 as amplifier-thermal-noise limited: the fundamental hardware floor.

\textbf{CH3} (rod piezo~A) shows \SI{+39}{\deci\bel} excess in the \SIrange{5}{10}{\kilo\hertz} band (\SI{68}{\milli\volt} RMS), from ambient coupling into the \SI{9.17}{\kilo\hertz} steel eigenmode even without active excitation. CH3 is paradoxically \emph{worse} than CH4 as a passive AE reference because structural resonance dominates.

\textbf{CH1} (soil piezo) carries \SI{137}{\milli\volt} total RMS, of which \SI{129}{\milli\volt} are in the \SIrange{0}{5}{\kilo\hertz} band. Contributing mechanisms: (i)~LM358 $1/f$ noise (dominant below \SI{1}{\kilo\hertz}); (ii)~capacitive mains coupling through soil; (iii)~ambient vibrations via soil matrix. AC coupling at the oscilloscope input ($f_{-3}<\SI{1}{\hertz}$) does not attenuate any of these.

%% ─────────────────────────────────────────────────────────────────────────────
\section{Plant AE Experiment}
%% ─────────────────────────────────────────────────────────────────────────────

\subsection{Protocol}

\begin{enumerate}[leftmargin=1.6em,itemsep=1pt,topsep=2pt]
  \item \textbf{Baseline} ($N_\text{pre}=10$ frames, \SI{0.6}{\second} each): continuous recording.
  \item \textbf{Watering}: single \SI{10}{\second} drip event via microcontroller.
  \item \textbf{Response} ($N_\text{post}=20$ frames): continuous recording immediately after.
\end{enumerate}

Total window: \SI{30}{\minute}. Plant: potted \textit{Euphorbia leuconeura}, $\approx18$ months.

\subsection{Pre--Post Watering Analysis}

A monotone upward trend is present within both the pre-watering sequence (Kendall $\tau=+0.69$, $p=0.005$) and the post-watering sequence ($\tau=+0.90$, $p<10^{-4}$) , consistent with thermal warming of the amplifier over time. The apparent $-\SI{2.9}{\milli\volt}$ drop at the watering boundary is a mechanical transient from physical disturbance, not an acoustic emission event.

The uncorrected Wilcoxon test (paired, pre vs.\ post) yields $p=0.002$ for CH1 RMS, nominally significant. This is a \emph{temporal confound}: the test compares pre-watering frames at the end of the warm-up curve (\SI{138.5}{\milli\volt}) against post-watering frames at the start of a new warm-up cycle (\SI{135.6}{\milli\volt}). After drift correction, comparing pre-steady-state (frames~0--4) against post-steady-state (frames~10--19), Mann--Whitney yields $p=0.86$, with means \SI{137.36}{\milli\volt} and \SI{137.42}{\milli\volt} respectively, a difference of \SI{0.06}{\milli\volt}.

Band-level FDR analysis nominally flags eight bands, but their directions are inconsistent (3 positive, 5 negative) and effect sizes negligible ($\max|\Delta E|=\SI{3.0}{\deci\bel}$). For comparison, the positive control in Section~\ref{sec:ref} shows \SI{+52.6}{\deci\bel} sensitivity. The two most significant bands (\SIrange{50}{55}{\kilo\hertz}) fall in the amplifier roll-off zone ($\gammasq=0.06$) and are uninformative.

\textbf{Result:} No plant AE signal is detectable. The \SI{137}{\milli\volt} noise floor on CH1 exceeds any credible plant AE amplitude by at least one order of magnitude.

\subsection{Continuous Monitoring}

A 51-hour continuous session logged 817 persistent cross-channel peak events. The dominant peak at \SI{1.5}{\kilo\hertz} (30th harmonic of \SI{50}{\hertz} mains) appears on all three channels simultaneously in 62\% of events, i.e.\ a common environmental source. Only 40 events (4.9\%) appear on CH1 but not CH3, the subset most likely to be soil-mediated; they cluster near \SI{6}{\kilo\hertz} with no biological timing pattern, consistent with the null watering result.

%% ─────────────────────────────────────────────────────────────────────────────
\section{MEMS Characterisation}
\label{sec:mems}
%% ─────────────────────────────────────────────────────────────────────────────

\subsection{Frequency Calibration and Noise Floor}
\label{sec:memscal}

Before the passive campaigns, the two MEMS channels (CH3, CH4) were characterised directly from raw data with the same estimators as the piezo chain (Sec.~\ref{sec:hw}). A Pico sweep (\SIrange{5}{100}{\kilo\hertz}, \SI{5}{\kilo\hertz} steps, 80 valid frames) drives the coherence, amplitude and MLE time-delay analysis (Fig.~\ref{fig:memscal}a,b), recomputed from the raw complex FFT bins; a passive deep-memory run gives the noise floor (Fig.~\ref{fig:memscal}c).

The calibrated band ($\gammasq\geq0.9$ and SNR${}\geq3$ on both channels) again spans \SIrange{5}{35}{\kilo\hertz}, bounded by the LM roll-off, with the strongest transfer at \SI{10}{\kilo\hertz} and \SI{30}{\kilo\hertz}. CH3 and CH4 track each other closely; across the reliable bands the sweep MLE time-delay stays within about \SI{1}{\micro\second} (excluding the low-SNR \SI{5}{\kilo\hertz} point), consistent with near-equidistant placement.


\begin{figure*}[!tb]
  \centering
  \includegraphics[width=0.82\textwidth]{figures/fig_mems_phase3.png}
  \caption{MEMS frequency calibration and noise floor. (a)~CH3--CH4 coherence $\gammasq$ per band; green shading marks the calibrated \SIrange{5}{35}{\kilo\hertz} band. (b)~Recorded RMS transfer amplitude of both MEMS channels (peaks at \SI{10}{} and \SI{30}{\kilo\hertz}). (c)~Passive noise floor: the two MEMS channels (CH3 \SI{81}{\milli\volt}, CH4 \SI{75}{\milli\volt} total AC RMS) sit above the soil piezo (CH1 \SI{39}{\milli\volt}) and are dominated by the common-mode comb in the \SIrange{5}{10}{\kilo\hertz} band.} 
\label{fig:memscal}
\end{figure*}

\subsection{Protocol and Decision Gates}

First an air reference was conducted just above the soil and then a follow-up with both microphones inserted into the soil at unequal distances from the presumed plant source region. Each run contains $N=20$ synchronised frames at \SI{500}{\kilo Sa\per s} and \num{300000} samples per channel.

The notebook method chain was applied without changing its decision logic: DC removal and linear detrending; clipping/plausibility rejection; Welch PSD ($\text{nperseg}=8192$); \SI{5}{\kilo\hertz} band energies and prominence peaks; robust Hilbert-envelope transient screening in \SIrange{20}{100}{\kilo\hertz}; CH3--CH4 coherence; narrowband phase MLE with GCC-PHAT cross-check; Kendall trend tests; and first-half versus second-half Mann--Whitney tests with Benjamini--Hochberg FDR. A credible plant event had to be transient rather than periodic, exceed the passive baseline, appear plausibly on both sensors, and exhibit a non-zero delay consistent with the unequal geometry.

\subsection{Air Reference and Soil Results}

In air, median RMS was \SI{80.67}{\milli\volt} on CH3 and \SI{79.38}{\milli\volt} on CH4 (CH4/CH3 $=-\SI{0.25}{\deci\bel}$). Seventeen of 20 frames lay within $\pm\SI{1}{\micro\second}$; after rejecting three low-confidence boundary solutions, the GCC-PHAT median was \SI{0.000}{\micro\second}. Phase candidates occurred mostly at the upper analysis boundary (\SIrange{90}{100}{\kilo\hertz}) and had \SIrange{196}{322}{\micro\second} uncertainty, far exceeding their fitted delays.

Soil insertion increased the median RMS to \SI{306.79}{\milli\volt} (CH3) and \SI{298.71}{\milli\volt} (CH4), while the relative level remained close (CH4/CH3 $=-\SI{0.23}{\deci\bel}$). The spectrum is a dense line comb shared by both channels; its strongest line at \SI{32.593}{\kilo\hertz} has coherence $\gammasq=0.997$. Mean \SIrange{20}{100}{\kilo\hertz} coherence is 0.603, and 14/16 five-kilohertz bands exceed 0.7. Nevertheless, all 20 GCC-PHAT estimates are exactly \SI{0.000}{\micro\second}; phase-MLE delays span only $-\SI{0.017}{\micro\second}$ to $+\SI{0.065}{\micro\second}$. This near-zero delay is inconsistent with using the unequal sensor distances to localise a soil-borne source and instead indicates common electrical or environmental coupling.

Transient screening produces approximately 544 threshold crossings per channel in every \SI{0.6}{\second} frame, with 10,392 CH3 candidates matched on CH4 within \SI{200}{\micro\second}. Their periodicity rejects them as AE clicks. Frame~13 contains the only large channel-specific excursion: CH3 steps from approximately $+\SI{1.4}{\volt}$ to $-\SI{2.54}{\volt}$ and remains at the rail while CH4 stays normal. The discontinuity and sustained plateau identify saturation/overflow, not an elastic transient. No first--second-half band comparison survives FDR correction (0/32); RMS trends are non-significant on CH3 (Kendall $\tau=-0.253$, $p=0.128$) and nominal on CH4 ($\tau=-0.358$, $p=0.028$) but unsupported by any FDR-significant spectral change.

\textbf{Result:} both MEMS channels respond and have similar average sensitivity, but neither the air nor soil run provides evidence of plant AE. Soil insertion mainly amplifies a coherent periodic interference field.



%% ─────────────────────────────────────────────────────────────────────────────
\section{Positive Control: Reference Channel Experiment}
\label{sec:ref}
%% ─────────────────────────────────────────────────────────────────────────────

To verify pipeline integrity, a reference experiment ($N=40$ frames) was conducted with the Pico sweep driving CH1 while CH2 captured the same signal and CH3 served as a fully passive reference (Fig.~\ref{fig:reference}).

CH1--CH2 multi-frame coherence at \SI{290}{\hertz} reaches $\gammasq=0.928$, versus $\gammasq\approx0.01$ for CH1--CH3, a \SI{+52.6}{\deci\bel} sensitivity advantage. Wilcoxon signed-rank on per-frame RMS (CH1 vs.\ CH3) gives $p<10^{-12}$. This confirms that the statistical pipeline detects genuine correlated signals when present, and that the negative watering result reflects a true absence of detectable AE above the noise floor.

\begin{figure*}[!tb]
  \centering
  \includegraphics[width=0.82\textwidth]{figures/fig7_reference.png}
  \caption{Positive control: a known Pico signal is fed into CH1 ($N=40$ frames). (a)~Coherence between the driven pair (CH1--CH2) and the passive pair (CH1--CH3); the driven pair tracks the signal (peak $\gammasq=0.928$ at \SI{290}{\hertz}), the passive pair stays near zero. (b)~Per-frame RMS: CH1 \SI{210}{\milli\volt} (driven), CH2 \SI{54}{\milli\volt}, CH3 \SI{0.45}{\milli\volt}. The \SI{52.6}{\deci\bel} gap shows the pipeline cleanly separates a real signal from a silent channel.}
  \label{fig:reference}
\end{figure*}

%% ─────────────────────────────────────────────────────────────────────────────

\begin{table}[!h]
\centering
\small
\caption{Noise floor and estimated SNR. Thermal limit: $e_n\sqrt{B}$,
         $e_n=\SI{40}{\nano\volt\per\sqrt{\hertz}}$. AE range from \citet{Cochard1992}.}
\label{tab:snr}
\begin{tabular}{lrll}
\toprule
Ch. & RMS (mV) & Band & Noise source \\
\midrule
CH4 &  25.7 & Wideband   & LM358 thermal \\
CH3 &  84.0 & 5--10\,kHz & Steel eigenmode \\
CH1 & 137.0 & 0--5\,kHz  & $1/f$ + mains \\
\bottomrule
\end{tabular}
\par\smallskip
{\footnotesize Plant AE (cavitation): $0.01$--$1$\,mV RMS \citep{Cochard1992};
SNR on CH1: $-43$ to $-23$\,dB (undetectable).}
\end{table}

%% ─────────────────────────────────────────────────────────────────────────────
\section{Discussion}
%% ─────────────────────────────────────────────────────────────────────────────

\subsection{Noise Floor as Primary Limitation}

CH1's \SI{137}{\milli\volt} floor (LM358 $1/f$ noise plus weak mechanical isolation) is $\approx\SI{80}{\deci\bel}$ above a low-noise pre-amplifier ($e_n\approx\SI{4}{\nano\volt/\sqrt{\hertz}}$, e.g.\ OPA2134). Reported plant AE amplitudes span tens of microvolts (single cavitation, \citealp{Cochard1992}) to a few millivolts, i.e.\ \SIrange{20}{40}{\deci\bel} below this floor. The LM358 also rolls off at \SI{35}{\kilo\hertz}, excluding the primary cavitation band (\SIrange{100}{500}{\kilo\hertz}); the usable \SIrange{5}{35}{\kilo\hertz} window, where slow hydraulic AE is reported \citep{Gagliano2012,Khait2023}, is noisy.

\subsection{Structural and Thermal Artefacts}

The passive \SI{+39}{\deci\bel} excess on CH3 at \SI{9.17}{\kilo\hertz} shows how a structural resonance amplifies ambient vibration straight into the band of interest; a damped reference mount or anechoic enclosure would remove it. Separately, the reproducible $+\SI{1.5}{\milli\volt}$ warm-up drift confounds sequential pre--post tests: future protocols should stabilise for \SI{30}{\minute}, interleave pre/post frames, and always report Kendall $\tau$ next to any Wilcoxon result on sequential data.

The MEMS follow-up reaches the same null decision by a different failure mode. Soil insertion raises both channels by approximately \SI{11.6}{\deci\bel} relative to the air reference, but the almost identical spectra, \SI{32.59}{\kilo\hertz} coherence of 0.997, and zero inter-channel delay show that the added energy is not a spatially propagating plant signal.

\subsection{Path Towards Detectability}

In order of expected impact: (1)~low-noise pre-amplifier ($\approx\SI{18}{\deci\bel}$); (2)~hardware \SI{1}{\kilo\hertz} high-pass filter to kill $1/f$ and mains ($\approx\SI{20}{\deci\bel}$ on CH1); (3)~resonant \SI{150}{\kilo\hertz} piezo with \SI{5}{\mega Sa\per s} ADC to reach the cavitation band; (4)~a vibration-isolated pedestal. Steps~1 and~2 together would bring CH1 below \SI{5}{\milli\volt} RMS, within range of strong AE events.

%% ─────────────────────────────────────────────────────────────────────────────
\section{Conclusion}
%% ─────────────────────────────────────────────────────────────────────────────

We characterised a piezoelectric AE system (calibrated band \SIrange{5}{35}{\kilo\hertz}, equidistant sensors) and evaluated it on \textit{Euphorbia leuconeura}. A controlled watering experiment yields no detectable AE response after thermal drift correction ($p=0.86$); the nominally significant Wilcoxon results ($p=0.002$) are fully explained by thermal drift (Kendall $\tau=+0.69$). A positive control confirms the pipeline detects real signals (\SI{+52.6}{\deci\bel}, $p<10^{-12}$), so the negative result is genuine, not a pipeline failure.

The MEMS campaign independently finds no AE evidence: air and soil runs both yield zero median delay, while soil insertion raises the median level to approximately \SI{303}{\milli\volt} and exposes a coherent \SI{32.59}{\kilo\hertz} interference comb plus a CH3 saturation fault. Similar channel sensitivity therefore confirms that both MEMS operate, but does not establish biological origin.

The limiting factors remain the LM358 noise floor, common-mode coupling, insufficient upper bandwidth, and incomplete geometry metadata. A low-noise pre-amplifier plus a hardware high-pass filter should lower the piezo floor by $\approx\SI{38}{\deci\bel}$; future MEMS runs additionally require recorded sensor distances, galvanic/common-mode controls, and a time-locked positive impulse before passive candidates can be attributed to the plant.

\vspace{4pt}
\noindent\textbf{Data and code.} All measurement data, analysis notebooks, and this paper's figure generation script are in the project repository at \texttt{src/characterization/}. Historical figures are reproducible via \texttt{python3 paper/generate\_figures.py}; the MEMS frequency calibration and noise floor (Fig.~\ref{fig:memscal}) is reproduced from the raw \mbox{git-LFS} recordings via \texttt{python3 scripts/analyze\_mems\_phase3.py}, and further results and plots reside with runs.

%% ─────────────────────────────────────────────────────────────────────────────
\FloatBarrier  % alle offenen Floats vor der Literaturliste platzieren
\bibliographystyle{plainnat}
{\footnotesize
\setlength{\bibsep}{1pt plus 0.3ex}
\begin{thebibliography}{99}

\bibitem[Benjamini and Hochberg(1995)]{Benjamini1995}
Y.~Benjamini and Y.~Hochberg (1995).
Controlling the false discovery rate: a practical and powerful approach to multiple testing.
\textit{J.\ R.\ Stat.\ Soc.\ B}, 57(1):289--300.

\bibitem[Bonisoli et al.(2025)]{Bonisoli2025}
L.~Bonisoli, L.~Forti, and L.~Arru (2025).
Outdoor detection of plant ultrasonic emissions using a contactless microphone.
\textit{Sensors}.

\bibitem[Cochard et al.(1992)]{Cochard1992}
H.~Cochard, P.~Cruiziat, and M.\,T.~Tyree (1992).
Use of positive pressures to establish vulnerability curves.
\textit{Plant Physiol.}, 100(1):205--209.

\bibitem[De Roo et al.(2016)]{DeRoo2016}
L.~De~Roo, L.\,L.~Vergeynst, N.\,J.\,F.~De~Baerdemaeker, and K.~Steppe (2016).
Acoustic emissions to measure drought-induced cavitation in plants.
\textit{Appl.\ Sci.}, 6(3):71.

\bibitem[Gagliano et al.(2012)]{Gagliano2012}
M.~Gagliano, S.~Mancuso, and D.~Robert (2012).
Towards understanding plant bioacoustics.
\textit{Trends Plant Sci.}, 17(6):323--325.

\bibitem[Khait et al.(2023)]{Khait2023}
I.~Khait et al.\ (2023).
Sounds emitted by plants under stress are airborne and informative.
\textit{Cell}, 186(7):1328--1336.

\bibitem[Kir\'{a}ly et al.(2025)]{Kiraly2025}
A.~Kir\'{a}ly, D.~Farkas, and J.~Dobr\'{a}nszki (2025).
Ultrasound in plant life and its application perspectives in horticulture and agriculture.
\textit{Plants}.

\bibitem[Klaminder et al.(2025)]{Klaminder2025}
J.~Klaminder, M.\,M.~Rollo~Jr, T.~N\"asholm, and J.~Viklund (2025).
Ultrasonic acoustic emissions as indicators of tree drought stress in outdoor forest settings.
\textit{For.\ Ecol.\ Manag.}

\bibitem[Knapp and Carter(1976)]{Knapp1976}
C.\,H.~Knapp and G.\,C.~Carter (1976).
The generalized correlation method for estimation of time delay.
\textit{IEEE Trans.\ Acoust.\ Speech Signal Process.}, 24(4):320--327.

\bibitem[Milburn and Johnson(1966)]{Milburn1966}
J.\,A.~Milburn and R.\,P.\,C.~Johnson (1966).
The conduction of sap. II. Detection of vibrations produced by sap cavitation in \textit{Ricinus} xylem.
\textit{Planta}, 69(1):43--52.

\bibitem[Nardini et al.(2015)]{Nardini2015}
A.~Nardini, M.~Battistuzzo, and T.~Savi (2013).
Shoot desiccation and hydraulic failure in temperate woody angiosperms during an extreme summer drought.
\textit{New Phytol.}, 200(1):322--329.

\bibitem[Oleti\'{c} et al.(2020)]{Oletic2020}
D.~Oleti\'{c}, S.~Rosner, M.~Zovko, and V.~Bilas (2020).
Time-frequency features of grapevine's xylem acoustic emissions for detection of drought stress.
\textit{Comput.\ Electron.\ Agric.}, 178:105787.

\bibitem[Oleti\'{c} et al.(2023)]{Oletic2023}
D.~Oleti\'{c}, S.~Rosner, and V.~Bilas (2023).
Field-experiences of tracking plant's xylem embolism formation with embedded acoustic emission sensors.
\textit{e-J.\ Nondestruct.\ Test.}, 28(3).

\bibitem[Saha and Rastogi(2026)]{Saha2026}
B.~Saha and A.~Rastogi (2026).
Plant acoustic emission as early stress signals: towards remote integrated monitoring for sustainable agriculture.
\textit{Eur.\ J.\ Agron.}

\bibitem[Simbeye et al.(2023)]{Simbeye2023}
D.\,S.~Simbeye, M.\,E.~Mkiramweni, and B.~Karaman (2023).
Plant water stress monitoring and control system.
\textit{Smart Agric.\ Technol.}, 5:100196.

\bibitem[Tyree and Dixon(1983)]{Tyree1983}
M.\,T.~Tyree and M.\,A.~Dixon (1983).
Cavitation events in \textit{Thuja occidentalis} L.?
Ultrasonic acoustic emissions from the sapwood can be measured.
\textit{Plant Physiol.}, 72(4):1094--1099.

\bibitem[Zweifel and Haesler(2001)]{Zweifel2012}
R.~Zweifel and R.~H\"asler (2001).
Dynamics of water storage in mature subalpine \textit{Picea abies}.
\textit{Tree Physiol.}, 21(9):561--569.

\end{thebibliography}
}

\end{document}
